\appendix
\section{附录：参数取值与求解步骤}

\subsection{题面给定常量与成本函数}
问题三的算力构成与质量成本参数按题面原样载入，列于表~\ref{tab:appendix-params}。

\begin{table}[!htbp]
  \centering
  \songti\zihao{-4}
  \caption{题面给定常量与三型质量成本参数}
  \label{tab:appendix-params}
  \begin{tabular}{lll}
    \toprule
    物理量 & 取值 & 说明 \\
    \midrule
    注意力系数 $\eta$ & $2\times10^{-4}$ & 作用于三者乘积 $N\Dt\Lctx$ \\
    Chinchilla 系数 & $6$ & 每参数每 token 的训练算力 \\
    注意力临界上下文 & $3\times10^{4}$ & 由 $6/\eta$ 解析得到 \\
    指数型 $g(Q)$ & $10^{7}\,e^{6Q}$ & 高质量极昂贵 \\
    幂型 $g(Q)$ & $5\times10^{9}\,Q^{4}$ & 高质量成本居中 \\
    对数型 $g(Q)$ & $2\times10^{9}\ln(1+10Q)$ & 高质量相对可得 \\
    三档预算 $C$ & $10^{19},10^{22},10^{24}$ & FLOPs \\
    基线质量 $Q_0$ & $0.48788$ & 问题一七域质量中位数 \\
    \bottomrule
  \end{tabular}
\end{table}

三型质量成本函数在中高预算下给出相近的最优质量，但在低预算档因高质量的经济性不同而分化，这是问题三比较成本函数选择影响的依据（表~\ref{tab:appendix-params}）。

\subsection{优化求解步骤}
问题三先利用预算活跃条件消去 $D$，在每个固定质量下求唯一的 $\ln N$ 驻点，再以质量网格、局部细化和端点比较选取候选。所有结果复算预算以及质量单边条件；另用 $513$ 点加密网格及 $20$ 起点 SLSQP 核验九组主预算。连续预算扫描输出份额导数和质量边界阈值。求解器一致性与局部残差不构成全局最优证明。

\subsection{关于人工智能工具的使用说明}
本文使用通用大语言模型辅助代码检查与修改、公式推导核验、图表生成、文字润色及 LaTeX 排版。报告数值以随文程序对题目附件的计算及已声明公式的复算为依据；模型假设、经验观测和情景外推分别标明，工具输出通过对应的数值检查后写入正文。
